\section{Varying density-matrix references}
\label{sec:ldp_varying_reference}

The reference state can vary with the ring and the site and can be singular.
We keep a fixed square bond space of dimension $D$. All reference matrices
$\sigma$ are positive semidefinite with trace one, and
$R_\sigma(X)=\operatorname{Tr}(X)\sigma$.
The ordered exponential estimates below are sufficient hypotheses; the
qualitative convergence definition in the inhomogeneous paragraph
of~\cite{Malz2024Preparation} does not supply them.
Rectangular bond spaces remain outside this statement
(\cite{gap:mswc24_inhomogeneous_scope}).

\begin{theorem}[Ordered perturbations of varying references]
    \label{thm:ldp_varying_reference_products}
    \lean{prod_range_sub_prod_range_eq_sum_mul_sub_mul,
        norm_prod_range_sub_prod_range_le_of_forall_norm_prod_le}
    \leanok
    Let $X_j,Y_j$ lie in a normed ring. Suppose every contiguous product of
    the $Y_j$, including the empty product, has norm at most $c$, and
    $\|X_j-Y_j\|\le\eta$. Then
    \begin{align}
        \left\|\prod_{j=0}^{M-1}X_j-\prod_{j=0}^{M-1}Y_j\right\|
        \le c\bigl((1+c\eta)^M-1\bigr).
    \end{align}
\end{theorem}
\begin{proof}\leanok
    Telescope the difference by taking an initial segment from $X$, one
    difference $X_j-Y_j$, and the remaining segment from $Y$. Bound each
    initial segment recursively by its reference product plus its error.
    The resulting geometric sum gives the assertion.
\end{proof}

\begin{theorem}[The reference on an outgoing pair]
    \label{thm:ldp_varying_reference_pairs}
    \lean{MPSPreparation.mpvFamily_fixedPointTensor,
        MPSPreparation.sum_mpvFamily_mul_star_mpvFamily}
    \leanok
    Write $F_j^{l,r}=\sqrt{\sigma_j}\,|l\rangle\langle r|$.
    For a periodic family of $M\ge1$ such tensors,
    \begin{align}
        \operatorname{Tr}(F_0^{l_0,r_0}\cdots F_{M-1}^{l_{M-1},r_{M-1}})
        =\prod_{j=0}^{M-1}(\sqrt{\sigma_{j+1}})_{r_j,l_{j+1}},
    \end{align}
    with cyclic subscripts. In particular the pair leaving block $j$
    uses the reference of block $j+1$. The overlap of two periodic tensor
    families is the trace of their ordered mixed-transfer product.
\end{theorem}
\begin{proof}\leanok
    Expand the cyclic virtual contraction in matrix units. Each right
    index fixes the incoming virtual index of the next tensor. Expanding
    both layers in the same way proves the mixed-transfer identity.
\end{proof}

The following two-block equality displays all four physical indices.
The lower path on the left closes only the virtual ring. On the right
there is no contraction between the two pair matrices: their coefficients
multiply. The matrix marked $\sqrt{\sigma_1}$ has row $r_0$ and column $l_1$;
the other has row $r_1$ and column $l_0$. Each pair has squared norm
$\operatorname{Tr}(\sigma_j)=1$, even when $\sigma_j$ is singular.
% TENKZ-ORDERED-MIXING-VARYING-BEGIN
% Formula: Tr(F0^(l0,r0) F1^(l1,r1)) = sqrt(sigma1)_(r0,l1) sqrt(sigma0)_(r1,l0).
% Boundary: l0,r0,l1,r1, each a physical index of dimension D.
\begin{center}
    \tenkzkernel
    \begin{tenkzequation}
        \begin{tenkz}[rows={wire,wire}, cols=6, bonds=none]
            \tn[at=(1,2), name=vpzero, skin=box, wide=2,
                ports={180:virtual, 0:virtual,
                        90@1:physical:$l_0$, 90@2:physical:$r_0$}]{F_0}
            \tn[at=(1,4), name=vpone, skin=box, wide=2,
                ports={180:virtual, 0:virtual,
                        90@1:physical:$l_1$, 90@2:physical:$r_1$}]{F_1}
            \tn[at=(1,1), name=vleft, skin=none, ports={0:virtual, 270:virtual}]{}
            \tn[at=(2,1), name=vleftbottom, skin=none, ports={90:virtual, 0:virtual}]{}
            \tn[at=(2,6), name=vrightbottom, skin=none, ports={180:virtual, 90:virtual}]{}
            \tn[at=(1,6), name=vright, skin=none, ports={180:virtual, 270:virtual}]{}
            \tnwire{vpzero.0}{vpone.180}
            \tnwire{vpzero.180}{vleft.0}
            \tnwire{vleft.270}{vleftbottom.90}
            \tnwire{vleftbottom.0}{vrightbottom.180}
            \tnwire{vrightbottom.90}{vright.270}
            \tnwire{vright.180}{vpone.0}
        \end{tenkz}
        $=$
        \begin{tenkz}[rows={wire}, cols=4, bonds=none]
            \tn[at=(1,1), name=vpairone, skin=box, wide=2,
                ports={90@1:physical:$r_0$, 90@2:physical:$l_1$}]{\sqrt{\sigma_1}}
            \tn[at=(1,3), name=vpairzero, skin=box, wide=2,
                ports={90@1:physical:$r_1$, 90@2:physical:$l_0$}]{\sqrt{\sigma_0}}
        \end{tenkz}
    \end{tenkzequation}
\end{center}
% TENKZ-ORDERED-MIXING-VARYING-END

\begin{theorem}[Uniform error with singular varying references]
    \label{thm:ldp_varying_reference_error}
    \lean{MPSPreparation.exists_norm_polarPos_sub_le_sqrt_transferMatrix,
        MPSPreparation.exists_norm_trace_prod_polarPos_sub_one_le_varying_reference,
        MPSPreparation.exists_abs_norm_chainState_sq_sub_one_le_uniform_reference,
        MPSPreparation.exists_one_sub_norm_inner_chainPosState_le_varying_reference}
    \leanok
    There are constants depending only on $D$ such that, if each actual
    block transfer matrix differs from its reset reference by at most
    $\delta$, its positive polar factor differs from
    $(\sqrt{\sigma_j})^\mathsf T\otimes I$ by at most $K_D\sqrt\delta$.
    The complex overlap of the positive-part state with the shifted pair
    family differs from one by at most
    $C_DM\sqrt\delta\exp(C_DM\sqrt\delta)$.
    If the whole-ring transfer matrix also differs from a density reset
    by at most $\delta$, the normalized pair error is at most
    $C'_DM\sqrt\delta$.
\end{theorem}
\begin{proof}\leanok
    \uses{thm:ldp_uniform_gram_transfer,thm:ldp_sqrt_holder,thm:ldp_varying_reference_products,thm:ldp_varying_reference_pairs,thm:ldp_inhomogeneous_normalized_pairs,thm:ldp_normalization_triangle,thm:ldp_chain_block_iso_inner}
    Apply the global square-root H\"older bound to the reshuffled Gram
    matrices. Density matrices form a compact set, so the reset matrices
    and the linear maps defining mixed transfers have uniform bounds.
    Since $R_\sigma R_\tau=R_\sigma$, every nonempty reference product has
    the same bound as one reset matrix and has trace one. Apply ordered
    telescoping, then use the whole-ring trace for the target norm.
    For $M\sqrt\delta\le1$ one has $\delta\le\sqrt\delta$ and the
    exponential is uniformly bounded. For the complementary range the
    normalized error is at most one.
\end{proof}

\begin{theorem}[Eventual preparation from varying-reference mixing]
    \label{thm:ldp_varying_reference_preparation}
    \lean{MPSPreparation.exists_isPreparedInDepth_inhomogeneous_le_log_eventually_of_varying_reference}
    \leanok
    Fix $d,D$ with $d>0$, and $K,r>0$. For every positive ring length let
    $A_{N,i}$ be actual tensors and $\sigma_{N,i}$ density matrices on
    the common bond space. Suppose every positive nonwrapping interval
    starting at $a$ and of length $\ell$ satisfies
    \begin{align}
        \|T_{A_{N,a}\cdots A_{N,a+\ell-1}}-R_{\sigma_{N,a}}\|
        \le K e^{-r\ell}.
    \end{align}
    There exist a uniform cutoff $N_0$ and $C$ such that the normalized
    original state at every $N\ge N_0$ and $0<\varepsilon\le1$ has a
    physical unit-state approximation of error at most $\varepsilon$
    prepared in depth $C\log(N/\varepsilon)$.
\end{theorem}
\begin{proof}\leanok
    \uses{thm:ldp_varying_reference_error,thm:ldp_inhomogeneous_uniform_rate_from,thm:ldp_chain_transfer_product,lem:ldp_log_threshold_error}
    The whole-ring estimate bounds the squared-norm deviation from one
    by $K_nK e^{-rN}$, uniformly over the density references. Choose the
    cutoff so this is at most $1/2$. On blocks of lengths between $q$ and
    $2q$ the preceding theorem gives error
    $C_D\sqrt K\,N e^{-rq/2}$. Apply the uniform-rate preparation theorem
    with exponent $r/2$, including its exact short-ring branch.
\end{proof}

\paragraph{Discussion and outlook.}
The constants now remain controlled for a family whose reference spectra
approach zero as the ring grows. Exact reset tensors, including different
rank-one references on neighboring blocks, satisfy the product hypothesis
with zero transfer error. This separates the dimension required by the
physical circuit from the rank of a reference density matrix. Faithfulness
is useful for the sharper linear square-root perturbation estimate, but
is not required for logarithmic preparation under the stated product bound.

The factor two lost in the exponential rate is a consequence of using a
global H\"older estimate, not a lower bound on preparation complexity.
A sharper estimate on supported sectors could improve it. Extending the
argument to rectangular bonds requires compatible corner-supported reset
maps and the corresponding pair supports. Deriving the uniform ordered
rate from a more intrinsic physical condition remains distinct from these
norm and normalization estimates.
